% Mid-year project progress report for CS 510 Intro to Privacy
% Steve Willoughby & Jacob Olinger
% 
% VLDB template version of 2020-08-03 enhances the ACM template, version 1.7.0:
% https://www.acm.org/publications/proceedings-template
% The ACM Latex guide provides further information about the ACM template
\documentclass[sigconf, nonacm]{acmart}
\usepackage{xcolor}
\definecolor{ontrack}{HTML}{4A6741}
\definecolor{behind}{HTML}{8F1402}

%% The following content must be adapted for the final version
% paper-specific
\newcommand\vldbdoi{XX.XX/XXX.XX}
\newcommand\vldbpages{XXX-XXX}
% issue-specific
\newcommand\vldbvolume{14}
\newcommand\vldbissue{1}
\newcommand\vldbyear{2020}
% should be fine as it is
\newcommand\vldbauthors{\authors}
\newcommand\vldbtitle{\shorttitle} 
% leave empty if no availability url should be set
\newcommand\vldbavailabilityurl{URL_TO_YOUR_ARTIFACTS}
% whether page numbers should be shown or not, use 'plain' for review versions, 'empty' for camera ready
\newcommand\vldbpagestyle{plain} 

\begin{document}
\title{IoT Access Control}
\subtitle{Mid-Quarter Progress Report}

%%
%% The "author" command and its associated commands are used to define the authors and their affiliations.
\author{Steve Willoughby}
\affiliation{%
  \institution{Portland State University}
  \streetaddress{P.O. Box 751}
  \city{Portland}
  \state{Oregon}
  \country{USA}
  \postcode{97207-0751}
}
\email{slw3@pdx.edu}
\author{Jacob Olinger}
\affiliation{%
  \institution{Portland State University}
  \streetaddress{P.O. Box 751}
  \city{Portland}
  \state{Oregon}
  \country{USA}
  \postcode{97207-0751}
}
\email{jacoboli@pdx.edu}
\maketitle
\section{Abstract}
Internet of Things (IoT) systems include a multitude of sensors which ``publish''
data streams to another multitude of ``subscribers''. Between them sits a ``broker''
which routes these published messages to all the subscribers interested in them.
However, there is a need for more robust access control to govern who can
subscribe to potentially sensitive data, and under what conditions they may do so.
Complicating things further, even an authorized subscriber may only have a valid
need to access sensor data at certain times, in certain locations or under certain
circumstances.

\section{Introduction}
\section{Related Work}
{\color{red}I have a bunch of content to paste in here from my thesis work\dots}
\section{Project Description}
Our project seeks to add fine-grained, high-performance access control to MQTT
pub/sub message broker systems for IoT networks.

It differs from the current state of the art in that most current offerings only go as far as basic
ACL rulesets.

\section{Timeline}
Our original estimated timeline was:
\begin{description}
	\item[\color{ontrack}15 Apr] Proposal Presentation.
	\item[\color{ontrack}19 Apr] Project scope and plan finalized following feedback.
	\item[\color{ontrack}26 Apr] Design for experimental simulated IoT network.
	\item[\color{ontrack}03 May] Implementation of IoT clients and brokers.
	\item[\color{behind}10 May] Control runs completed, Mid-Quarter Report.
	\item[\color{behind}17 May] Add access controls to broker, 2nd experiments.
	\item[24 May] Additional experiments and scale tests.
	\item[31 May] Complete data analysis.
	\item[05 Jun] Final Report ready.
\end{description}
\subsection{Progress to Date}
As suggested by the color-coded dates above, we are about a week behind schedule due to the time
required to search for suitable IoT client programs to use for publishers and subscribers and for
the time to fully settle on the access control policy model to be used in the broker modifications
we will be employing. 

We have obtained the clients and have an installed broker instance, and are just about ready to begin the
control runs now.

\section{Results and Analysis}
We will set up simulated traffic and measure the performance impact of our new
controls on the base system

We will run a number of test cases to demonstrate that the access controls are effective
at limiting access to only authorized subscribers who meet the policy restrictions within
the given domains.

\subsection{Experimental Setup}
We are working in a compartmentalized environment, which allows us to develop the server (MQTT broker) and clients (simulated publishing and subscribing devices) in relative isolation. The clients
are in Docker containers, allowing them to be instantiated identically as many times as needed in the experimental environment, each running inside its isolated container but able to send and receive network traffic with the server.

This setup allows each team member to separately develop the server and client components away from the actual test environment as well, allowing development and debugging in a way that does not interfere with running experiments.

The server used to perform all the experimental runs, and upon which our performance benchmarks is based, was an Intel\textregistered\ Core\texttrademark\ i9-10980XE (36 cores) running at 4.80\,GHz configured with 128\,Gb of memory.
It was running the Arch-based Garuda Linux distribution 6.13.7 Zen Linux kernel at the time of the experimental results reported here.
\section{Conclusions}
\section{Individual Contributions}
The efforts were divided between Steve and Jacob along the boundaries of broker and client systems,
with Steve handling the broker, including the setup of the Mochi MQTT broker and making the access
control enforcement modifications to it, and Jacob setting up the containerized instances of
the MQTT client programs acting as publishers and subscribers as well as the orchestration of
those clients' actions in sending their traffic as part of our experimental runs.
\section{References}
\end{document}
% 2-4 pp
%Introduction
%Describe the problem that you are addressing in this project and why it is interesting in the context of Privacy-aware Computing. Include examples or case studies to motivate the problem. Briefly describe what you hope to achieve at the completion of the project.
%	
%This criterion is linked to a Learning Outcome Related Work
%Overview of previous work related to project including research papers, projects, applications, or regulations. Importantly, how does your project relate to these existing approaches?
%
%Cite at least 3 other works.
%	
%This criterion is linked to a Learning Outcome Project Description
%A detailed description of the project including details of the specific privacy problem that you are addressing, definitions of terminology, details of the solution/approach, how the solution will improve privacy or privacy-awareness, how you are implementing the solution so far - details of tools/algorithms/frameworks that you are using, how you plan to evaluate your approach and what will be the metrics used. Include any architectural diagram or system diagrams of your approach. This will most likely be largest section of the report.
%	
%This criterion is linked to a Learning Outcome Timeline
%Include details of the tasks completed in the project and milestones achieved as well as the remaining tasks and how you plan to complete before the project presentation on March 10th. Optionally, include mention any technical or logistical challenges your team is facing and how you plan to address them.
%	
%This criterion is linked to a Learning Outcome Others
%Project title + Abstract
%Individual Contributions
%Results & Analysis
%This criterion is linked to a Learning Outcome References and Writing
%Correctly cited references
%Correctness of writing (Grammar, readability, spelling check)
